\subsection{Uncertainty budget and coverage}
\label{sec:coverage}

\begin{table}[t]
\centering
\caption{Standard uncertainty (mm) at the stump plane by input group (JCGM 100), root sum of squares within
a group: a representative delivery (120~km/h, 6~m, 60~Hz) and the median over \BudDistN{} deliveries.}
\label{tab:budget}
\scriptsize
\setlength{\tabcolsep}{3pt}
% generated by experiments/make_numbers.py; do not edit
\begin{tabular}{llcccc}
\toprule
& & \multicolumn{2}{c}{Representative} & \multicolumn{2}{c}{Median of 100} \\
\cmidrule(lr){3-4}\cmidrule(lr){5-6}
Input group & Type & $u(y)$ & $u(z)$ & $u(y)$ & $u(z)$ \\
\midrule
Calibration taps & B & 3.0 & 4.8 & 3.4 & 4.8 \\
Ball centre, bounce split & A, B & 7.3 & 35.2 & 15.6 & 25.1 \\
Frame timing & B & 3.0 & 16.5 & 9.5 & 14.9 \\
Scene dimensions & B & 1.4 & 12.6 & 4.1 & 15.3 \\
Lens distortion & B & 0.1 & 0.7 & 0.2 & 0.9 \\
Unmodelled flight & B & 38.5 & 50.8 & 37.9 & 37.7 \\
\midrule
Combined $u_c$ (JCGM 100) & & 39.4 & 65.3 & 46.9 & 53.8 \\
Monte Carlo $u$ (JCGM 101) & & 55.0 & 53.6 & & \\
Pipeline's own $\sigma$ & & 9.4 & 44.6 & & \\
\bottomrule
\end{tabular}

\end{table}

\begin{table}[t]
\centering
\caption{Coverage (\%) of the stated uncertainty on \CovBFitted{} World~B deliveries, which add every input
of Table~\ref{tab:budget}. The model term is cross-validated over five held-out folds.}
\label{tab:coverage}
\scriptsize
\setlength{\tabcolsep}{2pt}
% generated by experiments/make_numbers.py; do not edit
\begin{tabular}{lcccc}
\toprule
& \multicolumn{2}{c}{Lateral $y$} & \multicolumn{2}{c}{Vertical $z$} \\
\cmidrule(lr){2-3}\cmidrule(lr){4-5}
1-$\sigma$ used & $|e| \le \sigma$ & $|e| \le 1.96\sigma$ & $|e| \le \sigma$ & $|e| \le 1.96\sigma$ \\
\midrule
Exp.~1, pipeline $\sigma$ (before) & 51.6 & -- & 61.4 & -- \\
World B, fit only (as shipped) & 53.6 & 84.1 & 60.8 & 85.3 \\
World B, fit + calibration & 56.4 & 85.8 & 61.4 & 86.3 \\
World B, fit + calibration + model & 63.9 & 90.7 & 85.5 & 97.7 \\
\midrule
Nominal & 68.3 & 95.0 & 68.3 & 95.0 \\
\bottomrule
\end{tabular}

\end{table}

Each input in Table~\ref{tab:budget} was propagated through the calibration solve and the anchored
fit~\cite{jcgm100}, by numerical Jacobians for vector inputs (taps, ball centres, frame times) and by the RMS
output change over its PDF for scalar ones, so a bias the pipeline does not correct counts in full. The
measurement groups alone give \BudRepMeasY{} and \BudRepMeasZ~mm on the representative delivery, against a
pipeline sigma of \BudRepPipeSigmaY{} and \BudRepPipeSigmaZ~mm: the covariance is a fair measurement-noise
sigma. What it leaves out is the flight model, almost all of it unmodelled swing. A Monte Carlo propagation
of the same inputs through the full solver~\cite{jcgm101} ($M = \McM$) gives a vertical 95\,\% interval of
$[\McLowZ, \McHighZ]$~mm, shifted from the linear $[\McGumLowZ, \McGumHighZ]$~mm because swing of either sign
leaves the fit low. The linear budget fails the JCGM~101 validation by up to \McDHighZ~mm, so it ranks the
inputs and the Monte Carlo interval is the one to quote.

Table~\ref{tab:coverage} tests the stated sigma. Rolling shutter is not the cause of the shortfall: with no
readout skew and with 30~ms of it, lateral coverage is \SwRsZeroCovFitY\% and \SwRsThirtyCovFitY\%. The
missing term is model error, \ModYParamOne~mm laterally and \ModZParamOne~mm per metre of bounce-to-stump
distance vertically. Lateral coverage at 95\,\% is still \CovBFullYNinetyfive\%, with an empirical factor of
\KEmpFullYNinetyfive{} rather than 1.96, so the decision band takes its $k$ from risk and not from a nominal
Gaussian.
